\documentclass[11pt,a4paper]{article}
\usepackage[utf8]{inputenc}
\usepackage[T1]{fontenc}
\usepackage[spanish,provide=*]{babel}
\usepackage{geometry,booktabs,longtable,array,xcolor,listings,hyperref,titlesec}
\geometry{margin=2.25cm}
\definecolor{azulWCD}{RGB}{17,67,112}
\definecolor{tituloColor}{RGB}{0,55,110}
\definecolor{reglaColor}{RGB}{0,90,160}
\hypersetup{colorlinks=true,linkcolor=azulWCD,urlcolor=azulWCD}
\lstset{basicstyle=\ttfamily\small,breaklines=true,frame=single,columns=fullflexible}
\titleformat{\section}{\normalfont\Large\bfseries\color{azulWCD}}
  {\thesection.}{0.6em}{}
\titleformat{\subsection}{\normalfont\large\bfseries}
  {\thesubsection.}{0.5em}{}

\begin{document}
\pagestyle{empty}

%------------------------------------------------
% PORTADA
%------------------------------------------------

\begin{titlepage}
\begin{center}

\vspace*{0.5cm}
\textcolor{reglaColor}{\rule{\linewidth}{1.2pt}}

\vspace{1.3cm}
{\Large\scshape Universidad Industrial de Santander}\\[0.2cm]
{\large Escuela de Física --- Facultad de Ciencias}\\[0.2cm]
{\large Doctorado en Física}

\vspace{2.6cm}
{\Huge\bfseries\color{tituloColor} Implementación revisada\\[0.2cm]
del PMT en el WCD agua--NaCl}

\vspace{2.2cm}
{\large\textbf{Informe técnico}}

\vfill

{\large\textbf{Autor}}\\[0.15cm]
Propuesta técnica para MEIGA\\[0.1cm]
wcd-nacl-crns-simulations

\vspace{1.2cm}
\textcolor{reglaColor}{\rule{0.5\linewidth}{0.6pt}}

\vspace{1cm}
11 de septiembre de 2026

\vspace{1cm}
\textcolor{reglaColor}{\rule{\linewidth}{1.2pt}}

\end{center}
\end{titlepage}

\pagestyle{plain}

\section{Resultado}
Se revisó conjuntamente la definición óptica del medio sensible del PMT, su
geometría, el detector sensible, la curva de eficiencia y el registro de física
óptica. La implementación resultante conserva el modelo efectivo empleado por
MEIGA: el sólido elipsoidal del PMT actúa como fotocátodo absorbente, la frontera
con el agua o solución salina aplica las pérdidas de Fresnel mediante los índices
de refracción y \texttt{OptDevice} aplica la respuesta espectral. No se añadió
\texttt{EFFICIENCY} a una \texttt{G4OpticalSurface}, pues ello duplicaría la
eficiencia cuántica ya aplicada por \texttt{G4MPMTAction}.

\section{Problemas encontrados}
\begin{longtable}{>{\raggedright\arraybackslash}p{0.29\textwidth}
                  >{\raggedright\arraybackslash}p{0.61\textwidth}}
\toprule
Problema & Consecuencia \\
\midrule\endhead
\texttt{ProcessHits} se ejecutaba en cualquier paso dentro del PMT &
Un fotón no aceptado por la QE podía volver a sortearse en pasos posteriores. \\
No se terminaba explícitamente el track & Un fotón aceptado podía producir más
de un fotoelectrón lógico; uno rechazado continuaba hasta ser absorbido. \\
QE escalonada en intervalos de 50 nm & Introducía discontinuidades artificiales
y asignaba 6\% entre 650--700 nm después de 3\% entre 600--650 nm. \\
Rango 250--700 nm frente a la banda descrita de 300--650 nm & Permitía fotones
fuera de la respuesta nominal adoptada para el R5912. \\
Tres archivos de depuración escritos por fotón & Alto coste de E/S, archivos
muy grandes y riesgo de mezcla en ejecución multihilo. \\
Uso de \texttt{GetDetectorSimData()} sin identificador & Riesgo de asociar la
información al detector incorrecto cuando existe más de un detector. \\
Contadores sin inicialización explícita & Valores indeterminados antes del
primer incremento. \\
Comentarios describían Pyrex como vidrio real & En realidad
\texttt{ABSLENGTH}=0.5~$\mu$m lo convierte en un absorbente efectivo. \\
\bottomrule
\end{longtable}

\section{Cambios implementados}
\subsection{G4MPMTAction.cc y G4MPMTAction.h}
La respuesta sólo se evalúa cuando el prepaso tiene estado
\texttt{fGeomBoundary}, es decir, al entrar al volumen del PMT. Se descartan las
partículas no ópticas. Los fotones fuera de banda, no detectados o detectados
terminan con \texttt{fStopAndKill}; únicamente los aceptados se añaden a las
distribuciones temporales. Así, cada track puede generar como máximo un
fotoelectrón. Se inicializan y reinician contadores y vectores por evento, y se
consulta \texttt{GetDetectorSimData(fDetectorId)}.

\subsection{OptDevice.cc}
Se mantiene la información nominal ya presente en el proyecto, pero se expresa
como nodos de una curva continua:
\begin{center}
\begin{tabular}{rrrrrrrrr}
\toprule
$\lambda$ (nm)&300&350&400&450&500&550&600&650\\
QE&0.03&0.20&0.25&0.20&0.14&0.07&0.03&0.00\\
PDE (QE$\times0.70$)&0.021&0.140&0.175&0.140&0.098&0.049&0.021&0.000\\
\bottomrule
\end{tabular}
\end{center}
Entre nodos se usa interpolación lineal. La probabilidad empleada por el sorteo
es QE multiplicada por la eficiencia de colección del primer dínodo, fijada en
0.70 en el código preexistente. La probabilidad se limita al intervalo $[0,1]$.
Se eliminaron los archivos \texttt{eficiencia-limpia*.txt} generados por cada
fotón. El rango energético queda en 1.9077--4.1333 eV, equivalente a
650--300 nm.

Esta curva conserva y suaviza la parametrización previa; no debe presentarse
como una digitalización metrológica de una hoja de datos. Antes de resultados
finales debe contrastarse con la curva del número de serie/modelo usado o con
una medida experimental.

\subsection{Materials.cc}
La tabla del medio efectivo del PMT cubre la misma banda 300--650 nm. Se mantiene
$n=1.47$ y \texttt{ABSLENGTH}=0.0005 mm porque, en esta arquitectura, el sólido
no representa una ventana macroscópica: es un fotocátodo efectivo terminal.
Los comentarios ahora explican esta aproximación y advierten que no debe
añadirse una segunda \texttt{EFFICIENCY}.

\subsection{SaltyWCD.cc}
Se valida al construir la geometría que el PMT tenga \texttt{RINDEX} y
\texttt{ABSLENGTH}. Un error produce una excepción fatal identificada como
\texttt{WCD-OPT-003}. También se informa en el log que la respuesta combina
Fresnel, absorción terminal y QE de \texttt{OptDevice}.

\section{Dependencia con la salinidad}
La QE y la eficiencia de colección pertenecen al PMT y son comunes a los cuatro
medios. La salinidad influye antes de la detección mediante el \texttt{RINDEX}
del medio detector: cambia el ángulo de Cherenkov, el transporte y el coeficiente
de Fresnel en la frontera medio--PMT. Esta separación evita atribuir al PMT una
dependencia química que corresponde al medio.

\section{Archivos revisados sin cambios}
\texttt{OpticalPhysics.cc/.hh} registran Cherenkov, scintillation, absorption,
boundary y WLS. Rayleigh y MieHG están construidos pero su registro permanece
comentado. No se activaron en este parche porque alteraría todo el transporte,
no sólo el PMT, y los datos de dispersión salina están documentados como una
aproximación común. Deben evaluarse en un cambio independiente.

\texttt{G4WCDSteppingAction} registra producción e interacciones, pero no debe
volver a contabilizar la detección del PMT mientras ésta sea responsabilidad de
\texttt{G4MPMTAction}. \texttt{Detector} y \texttt{OptDevice.h} no requieren
cambios de interfaz para este modelo.

El archivo suministrado como \texttt{Materials(1).h} no es el encabezado de
\texttt{Materials}: su contenido corresponde a una implementación de
\texttt{SaltyWCD::BuildDetector}. Por ello no debe copiarse como
\texttt{Materials.h}. Este parche no añade miembros a \texttt{Materials} y puede
usar el encabezado real ya existente en el repositorio.

\section{Validación realizada y pendiente}
Se verificaron estáticamente monotonía del dominio espectral, límites de QE/PDE,
conversión energía--longitud de onda, ausencia de las escrituras de depuración,
y presencia de la condición de frontera y terminación del track. El entorno de
trabajo no contiene \texttt{geant4-config} ni el árbol completo de cabeceras de
MEIGA; por tanto no fue posible compilar ni ejecutar una corrida Geant4 aquí.

Antes de producción se requieren pruebas monoenergéticas, para cada uno de los
cuatro medios, con fotones dirigidos al PMT. Debe comprobarse:
\begin{enumerate}
\item que ningún \texttt{trackID} produce más de un PE;
\item que la fracción detectada, condicionada a entrar al PMT, converge a la PDE;
\item que la fracción que alcanza/entra cambia con el índice del medio;
\item que el balance entre reflexión, absorción y detección cierra;
\item que el resultado es reproducible con semilla fija y seguro en multihilo.
\end{enumerate}

\section{Limitaciones y siguiente nivel de fidelidad}
El modelo efectivo no separa vidrio, vacío y fotocátodo, ni implementa
reflectividad propia del fotocátodo o el 90\% de cobertura activa comentado en
versiones previas. Para esa fidelidad se requiere una cáscara elipsoidal de
vidrio con espesor real, una superficie interna \texttt{dielectric\_metal},
curvas medidas de $n(\lambda)$, transmisión, reflectividad y QE, y un único
mecanismo de registro del estado \texttt{Detection}. No debe añadirse esa capa
sin disponer de dichos datos, pues cambiaría la definición estadística de QE.

\section{Resumen listo para commit}
\begin{lstlisting}
fix(pmt): apply optical response once at photocathode entry

- gate PMT sensitive response on fGeomBoundary
- terminate accepted, rejected and out-of-range optical photons
- prevent duplicate photoelectrons from repeated Pyrex steps
- interpolate the nominal R5912 QE over 300--650 nm
- apply collection efficiency once and remove per-photon debug files
- align PMT material and OptDevice spectral ranges
- validate effective PMT RINDEX and ABSLENGTH during geometry build
- document the effective-absorber model and its limitations
\end{lstlisting}

\section{Referencias}
Geant4 Collaboration, \emph{Book for Application Developers}, procesos ópticos:
\url{https://geant4.web.cern.ch/documentation/pipelines/master/bfad_html/ForApplicationDevelopers/TrackingAndPhysics/physicsProcess.html}

Geant4 Collaboration, \texttt{G4OpBoundaryProcess}:
\url{https://github.com/Geant4/geant4/blob/master/source/processes/optical/src/G4OpBoundaryProcess.cc}

\end{document}
